% ============================================================================
%  figura3_tikz.tex  -  Diagrama de estados de la FSM del transmisor MAX7219
%  Pong en FPGA - Electrónica Digital 1 - UNAL
%  Daniel Felipe Castro González, Juan David Ortiz Muñoz, David Eduardo Chivata Arias
% ----------------------------------------------------------------------------
%  Máquina de Moore: las salidas CLK_LED y LOAD dependen solo del estado.
%  Entradas: T = se cumplió el tiempo del estado, U = es el último bit.
%  Preámbulo: \usetikzlibrary{automata,positioning,arrows.meta}
% ============================================================================
\begin{tikzpicture}[
    font=\footnotesize,
    >={Stealth[length=2.2mm]},
    every state/.style={draw, circle, minimum size=2.05cm, align=center, inner sep=1pt, line width=0.7pt},
    tr/.style={->, line width=0.7pt},
    lbl/.style={font=\scriptsize, fill=white, inner sep=1.5pt},
    node distance=3.4cm,
]
\node[state, initial, initial text={reinicio}, fill=orange!25] (carga) {\textbf{CARGA}\\000\\[1pt]{\scriptsize CLK=0, LOAD=0}};
\node[state, fill=cyan!15, right=of carga] (bajo) {\textbf{BAJO}\\001\\[1pt]{\scriptsize CLK=0, LOAD=0}};
\node[state, fill=blue!25, right=of bajo] (alto) {\textbf{ALTO}\\010\\[1pt]{\scriptsize CLK=1, LOAD=0}};
\node[state, fill=yellow!30, below=2.6cm of alto] (cierre) {\textbf{CIERRE}\\011\\[1pt]{\scriptsize CLK=0, LOAD=0}};
\node[state, fill=red!25, left=of cierre] (latch) {\textbf{LATCH}\\100\\[1pt]{\scriptsize CLK=0, LOAD=1}};
\node[state, fill=gray!20, left=of latch] (separa) {\textbf{SEPARA}\\101\\[1pt]{\scriptsize CLK=0, LOAD=0}};
\node[draw, dashed, rounded corners, align=center, above=0.9cm of carga, xshift=-2.4cm] (inv) {110, 111\\{\scriptsize (no usados)}};

\path[tr]
  (carga) edge node[lbl, above] {siempre} node[lbl, below] {\scriptsize sr $\leftarrow$ trama} (bajo)
  (bajo)  edge[bend left=18] node[lbl, above] {$T$} (alto)
  (alto)  edge[bend left=18] node[lbl, below] {$T\,\overline{U}$ \scriptsize(desplazar)} (bajo)
  (alto)  edge node[lbl, right] {$T\,U$} (cierre)
  (cierre) edge node[lbl, above] {$T$} (latch)
  (latch) edge node[lbl, above] {$T$} (separa)
  (separa) edge node[lbl, right] {$T$ \scriptsize(idx $\leftarrow$ siguiente)} (carga)
  (bajo)  edge[loop above] node[lbl] {$\overline{T}$} ()
  (alto)  edge[loop above] node[lbl] {$\overline{T}$} ()
  (cierre) edge[loop below] node[lbl] {$\overline{T}$} ()
  (latch) edge[loop below] node[lbl] {$\overline{T}$} ()
  (separa) edge[loop below] node[lbl] {$\overline{T}$} ();
\draw[tr, dashed] (inv) -- node[lbl, above right] {siempre} (carga);

\node[align=left, anchor=north west, font=\scriptsize] at ([xshift=1.0cm, yshift=-0.2cm]cierre.south east)
  {$T$: se cumplió el tiempo del estado\\(medio periodo de CLK\_LED;\\en LATCH, dos medios periodos)\\[3pt]
   $U$: es el último bit (bits $=1$)\\[3pt]
   Se repite 64 veces BAJO $\leftrightarrow$ ALTO\\(16 bits por chip $\times$ 4 chips)};
\end{tikzpicture}
